\documentclass{article}
\pdfpagewidth=8.5in
\pdfpageheight=11in

\usepackage{ijcai26}
\usepackage{times}
\usepackage{url}
\usepackage[hidelinks]{hyperref}
\usepackage[utf8]{inputenc}
\usepackage{graphicx}
\usepackage{amsmath}
\usepackage{booktabs}
\usepackage[brazilian]{babel}

\urlstyle{same}

% Portuguese abstract heading. The IJCAI style writes "Abstract" directly, so we
% redefine it. Babel already translates "References".
\renewenvironment{abstract}{\centerline{\Large\bf
Resumo}\vspace{0.5ex}\begin{quote}}{\par\end{quote}\vskip 1ex}

% Required by the IJCAI template. Leave this block untouched.
\pdfinfo{
/TemplateVersion (IJCAI.2026.0)
}

\title{Resolvendo o ARC-AGI-1 com LLMs: Fim a Fim e Geração de Código\\(Nível 2)}

\author{
    Rafaela Lorenzini
    \affiliations
    Especialização em Inteligência Artificial Avançada, UFRGS
    \emails
    rafaela.lorenzini@gmail.com
}

\begin{document}

\maketitle

\begin{abstract}
    O ARC-AGI-1 exige inferir uma regra a partir de poucos pares de grades e aplicá-la a uma
    entrada nova. Este trabalho compara duas formas de usar o Gemini 3.5 Flash Lite, com
    raciocínio alto, em 20 tarefas do benchmark: na abordagem fim a fim, o modelo escreve a
    grade de saída; na geração de código, ele escreve uma função Python executada pelo
    computador. A geração de código resolveu 18 das 20 tarefas e a abordagem fim a fim, 17,
    com custo de tokens e tempo semelhante. As duas abordagens erraram tarefas diferentes, o
    que sugere que elas se complementam.
\end{abstract}

\section{Introdução}

O ARC-AGI-1 avalia a capacidade de inferir regras abstratas a partir de poucos
exemplos~\cite{chollet2019arXiv}. Cada tarefa traz de dois a cinco pares de grades de entrada
e saída, com cores codificadas de 0 a 9, e uma ou mais entradas de teste. A resposta só conta
como certa quando a grade inteira coincide com a esperada. Cada tarefa usa uma regra própria,
como completar formas, contar objetos ou recortar uma região, e por isso não basta memorizar
padrões: o modelo precisa formular uma hipótese e aplicá-la sem erro em cada célula.

Neste trabalho, usamos o Gemini~\cite{geminiteam2025arXiv} de duas formas. Na abordagem
\emph{fim a fim}, o modelo lê os exemplos e escreve a grade de teste. Na \emph{geração de
código}, ele escreve um programa, e o computador executa esse programa. A segunda forma separa
a descoberta da regra da sua aplicação e permite conferir o programa nos exemplos de treino.

\section{Método}

O solver foi escrito em Python e roda com o comando \texttt{uv run solver/main.py}. Ele lê a
lista \texttt{tarefas.txt}, resolve cada tarefa nas duas abordagens e salva, para cada tarefa,
o prompt, a resposta bruta do modelo e um arquivo com o resultado, os tokens e o tempo. Ao fim
de cada abordagem, grava um resumo em JSON e em CSV. Há uma pausa de 2~s entre tarefas, e
erros temporários da API (429 ou 5xx) são repetidos até três vezes.

\paragraph{Prompt comum.} As duas abordagens começam pelo mesmo texto: uma instrução curta
sobre o problema, os pares de treino e as entradas de teste. Cada grade aparece como linhas de
números separados por espaço.

\paragraph{Fim a fim.} O prompt pede um objeto JSON com um resumo da regra em uma ou duas
frases e a lista de grades de saída, uma por entrada de teste. O solver extrai o JSON da
resposta e compara cada grade com a esperada. Quando a tarefa tem uma única entrada de teste,
o modelo às vezes envia a grade fora da lista; o solver aceita esse formato.

\paragraph{Geração de código.} O prompt pede uma função \texttt{transform(grid)}, que recebe
uma grade como lista de listas e devolve a grade de saída. O modelo pode usar a biblioteca
padrão, o NumPy e o SciPy. O solver extrai o bloco de código da resposta e executa a função
em um processo Python separado, com tempo limite de 10~s, em cada exemplo de treino e em cada
entrada de teste. O programa conta como certo no treino quando reproduz todas as saídas
conhecidas. A resposta ao teste é a saída do programa, mesmo quando ele erra o treino. Não há
segunda tentativa: cada tarefa recebe uma única chamada ao modelo.

\section{Experimentos}

Usamos o Gemini 3.5 Flash Lite (\texttt{gemini-3.5-flash-lite}) pela API do Google, com nível
de raciocínio alto nas duas abordagens. O conjunto tem as 20 tarefas do projeto: 10 do
conjunto de treino e 10 do conjunto de avaliação do ARC-AGI-1. Os tokens de entrada medem o
prompt; os de raciocínio, o texto que o modelo gera antes da resposta; e os de saída, a
resposta final. O tempo soma, por tarefa, a chamada à API e, na geração de código, a execução
do programa, sem as pausas entre tarefas. O solver foi construído com o Claude Code.

\begin{table}[h]
  \centering
  \small
  \setlength{\tabcolsep}{3pt}
  \begin{tabular}{lrrrrr}
    \toprule
    Abordagem   & Resolvidas & Entrada & Racioc. & Saída & Tempo (s) \\
    \midrule
    Fim a fim   & 17/20      & 37.276  & 198.615 & 5.903 & 570,0     \\
    Código      & 18/20      & 37.976  & 202.445 & 7.284 & 611,2     \\
    \bottomrule
  \end{tabular}
  \caption{Tarefas resolvidas, total de tokens e tempo por abordagem. Na geração de código,
  o programa acertou todos os exemplos de treino em 18 tarefas.}
  \label{tab:resultados}
\end{table}

\section{Resultados}

A Tabela~\ref{tab:resultados} mostra que a geração de código resolveu uma tarefa a mais que a
abordagem fim a fim, com custo próximo: cerca de 2\% a mais de tokens de raciocínio e 7\% a
mais de tempo. Os tokens de raciocínio dominam o custo nas duas abordagens e superam em mais
de cinco vezes os de entrada. Na geração de código, o programa acertou todos os exemplos de
treino em 18 tarefas, e todas elas também acertaram o teste. Assim, acertar o treino foi um
bom sinal de acerto no teste neste conjunto.

As abordagens erraram tarefas diferentes. A abordagem fim a fim errou c658a4bd, d47aa2ff e
f5aa3634; a geração de código errou 32597951 e f5aa3634. Juntas, elas cobririam 19 das 20
tarefas. Também observamos variação entre execuções: em outra rodada completa, a abordagem fim
a fim também resolveu 17 tarefas, mas não exatamente as mesmas.

\paragraph{Exemplo de acerto (3aa6fb7a).} A grade tem formas de cor 8 com três células em um
quadrado 2$\times$2, e a saída pinta de cor 1 a célula vazia de cada quadrado. As duas
abordagens acertaram. A abordagem fim a fim descreveu a regra como ``para cada grupo de três
células 8 em forma de L, colocar 1 no canto que falta''. O programa gerado usa a função
\texttt{label} do SciPy para achar cada componente da cor 8, calcula o retângulo que o
contém e pinta de 1 as células vazias desse retângulo. O programa acertou os dois exemplos de
treino e o teste.

\paragraph{Exemplo de erro (f5aa3634).} A entrada tem vários objetos pequenos, e a saída é o
recorte do objeto que aparece mais de uma vez. As duas abordagens erraram. Na abordagem fim a
fim, o modelo inferiu uma regra diferente (``o bloco indicado pelas formas marcadoras à
esquerda'') e devolveu uma grade 3$\times$3, enquanto a resposta tem 4$\times$3. Na geração de
código, o programa acertou dois dos três exemplos de treino e errou o teste. É a tarefa com
mais tokens de raciocínio na geração de código (21.829): o modelo gastou mais e não achou a
regra. Nesse caso, o erro no treino já indicava que o programa não era confiável.

\section{Conclusão}

Com o Gemini 3.5 Flash Lite e raciocínio alto, a geração de código foi a melhor abordagem:
resolveu 18 das 20 tarefas, contra 17 na abordagem fim a fim, a um custo semelhante em tokens
(cerca de 248 mil contra 242 mil no total) e em tempo. Além do acerto, a geração de código tem
uma vantagem prática: o resultado no treino indica quando a resposta é confiável. Como as
abordagens erram tarefas diferentes, um próximo passo natural é corrigir o programa com o erro
dos exemplos de treino (nível 3) ou combinar as duas respostas.

O solver foi construído com o agente de codificação Claude Code, que também ajudou a revisar
os resultados e a redigir este relatório; os números foram conferidos com os resumos da pasta
\texttt{resultados/}.

%% The file named.bst is a bibliography style file for BibTeX 0.99c
\bibliographystyle{named}
\bibliography{referencias}

\end{document}
